% Latex Document modified from Latex design of Brian R. Hall (available at www.brianrhall.net)

% Document settings
\documentclass[11pt]{article}
\usepackage[margin=1in]{geometry}
\usepackage[pdftex]{graphicx}
\usepackage{multirow}
\usepackage{setspace}
\pagestyle{plain}
\usepackage{color}
\setlength\parindent{0pt}
\makeatletter
\setlength{\@fptop}{0pt}
\makeatother

\begin{document}

% Course information
  \begin{center}

{\includegraphics[height=1.25in,width=1in]{wmchiffre1.jpg}} 

\LARGE Data Driven Decisionmaking (Fall 2015, INTR 490 (COLL 200))\\ \vspace{3mm}
\end{center}
\large \textbf{Schedule} \\
\normalsize  Monday, Wednesday, and Friday 11AM to 11:50AM \textit{Morton Hall 20} \\
\vspace{2mm}

\begin{table}[ht]
\begin{tabular}{l l}
\large \textbf{Instructors} & \\
\large Dr. \textbf{Dan} Runfola & Dr. Tyler Frazier \\
\large dsmillerrunfol@wm.edu  & tfrazier@aiddata.org \\
\large 427 Scottland Street & 427 Scottland Street \\
\large Office Hours: TBD & TBD \\
\large 757.221.1970 & 757.221.5396 \\
\end{tabular}
\end{table}


\large \textbf{Teaching Assistants} \\
\large Name: TBD  \\
\large Office: TBD \\
\large Office Hours: TBD \\
\large Number: TBD \\


\textit{Please let me know if you have any documented disabilities that may impact your performance in this class.}

\textbf {\large \\ Course Description:} For the increasing volume of information being produced to be useful in decision-making processes, it needs to be systematically organized and analyzed. This course will provide students with an opportunity to apply quantitative methods to a wide variety of real-world problems defined by decision makers from federal and international policy making groups. Course work will include applying a wide set of techniques (such as the analytic hierarchy process, ordered weighting averaging) which integrate human preferences and perception with quantitative information, with an emphasis on uncertainty. Students will learn about and consider the challenges associated with data reduction – how to balance between the limits of human perception, the value of additional information, and temporal constraints imposed by the decision making process.  \\
\\
\textbf {Prerequisite(s):} None.

\textbf {Credit Hours:} 3 \\

\textbf {\large Books and other things you'll need to buy:}\\ 

Books: None, though throughout the course I may assign readings from academic or popular journals – these will be made available through either the Swem library or online via blackboard.\\
Software: You must have a laptop capable of running Microsoft Excel (free versions, such as open or libre office, will not work for this course due to the complexity of some spreadsheets).


\vspace{8mm}

\textbf {\large Course Objectives:} 
\begin{enumerate} \itemsep-0.4em
  \item Provide students with a critical understanding of the decision making process and the use of data and intuition on short deadlines.
  \item Develop student’s ability to communicate findings, analysis, and visualization skills for future courses (and jobs).
  \item Expose students to real-world problems that are being engaged with by contemporary problem solvers and decision makers.
  \item Provide an opportunity to earn credit towards the COLL 200 requirement in the Natural World and Quantitative Reasoning domain (passing grade required).
\end{enumerate}
\vspace{8mm}

\textbf {\large Grade Distribution:} \\
\hspace*{40mm}
\begin{tabular}{ l l }
Lab Assignments & 50\% \\
Final Assignment & 30\%\\
Reading Assignments & 20\%\\ \\
\end{tabular} \\\\

\textbf {\large Letter Grade Distribution:} \\
\hspace*{40mm}
\begin{tabular}{ l l | l l }
\textgreater= 93.00 & A & 73.00 - 76.99 & C \\
90.00 - 92.99 & A-  & 70.00 - 72.99 & C- \\
87.00 - 89.99 & B+  & 67.00 - 69.99 & D+ \\
83.00 - 86.99 & B  & 63.00 - 66.99 & D \\
80.00 - 82.99 & B-  & 60.00 - 62.99 & D- \\
77.00 - 79.99 & C+  & \textless= 59.99 & F \\
\end{tabular} \\
.\\
% Course Policies. These are just examples, modify to your liking.
\textbf {\large Time Commitment:} Excelling in college level course work typically requires on average three to four hours per credit per week.  Since this is a four credit course, in addition to the time we meet as a class each week, you should expect to spend nine to twelve hours on average reading, writing, or otherwise preparing for this class on a weekly basis.\\

\textbf {\large Attendance:} This class does not have an attendance policy.  However, it will be difficult to learn enough to pass the class without regular participation, as the majority of course content relevant to tests and assignments will be covered in class.  Unannounced opportunities for extra credit will periodically be given in-class.\\

\textbf {\large Discussions:} Most lecture sessions will begin with a brief discussion of the assigned materials.  As such, each week students will be asked to come to class prepared with two to three bullet points for discussion - these can represent questions the material raised, commentary, or critiques.  These discussion points may be collected for extra credit.\\

\textbf {\large Classroom Behavior:} Please remain civil during discussions to promote the open exchange of ideas and foster a culture of open dialogue.  Please bear in mind that all students are entitled to their own opinion.  You are expected to listen attentively to each person speaking.  Please refrain from eating during class (and, if you must, make sure it isn't loud!).\\

\textbf {\large Teacher-student conferences:} Students performing at a C level or below are required to schedule a meeting with the instructor to discuss class performance.\\

\textbf {\large Late / Poor Performance Policy:} Assignments will not be accepted late, excepting in documented circumstances (i.e., an illness with a doctor's note).  \textit{Recognizing that it is entirely possible to underperform on key days, your lowest lab assignment grade will be dropped at the end of the semester.}\\

\textbf {\large Final Project:} The final project will be the culmination of your work in the second half of the semester, and will include both a written report and verbal presentation.  Final presentations will begin the final week of the course, and finish on the day of finals. \\

\vspace{4mm}
\textbf {\LARGE Do not cheat!} \\
.\hrulefill . \\
\textbf{Academic dishonesty is taken very seriously.  Make sure to cite all of your work, and do not turn in work that is not yours!  Cases of academic dishonesty will be evaluated and acted upon in accordance with William and Mary policies, which can be found at http://www.wm.edu/offices/deanofstudents/services/
studentconduct/} \\
.\hrulefill . \\
\vspace{8mm}

% Course Outline
\textbf {\large Course Outline}:

The course outline can be found on the following two pages.  The weekly content might change as it depends on the progress of the class.

\begin{table}[h!]
\small % The size of the table text can be changed depending on content. Remove if desired.
\begin{tabular}{ | c | c | }
\hline
\textbf{Week} & \textbf{Content} \\
\hline
Week 1 & \begin{minipage}{.85\textwidth}
\begin{itemize} \itemsep-0.4em
	\vspace{1mm}
	\item 8/26 - Lecture; 8/28 - Hands On Learning
	\item Data Applications in  Decisionmaking: A History I  \item Reading: \textit{A Primer for Policy Analysis, Ch. 3- Decision Making by  Organizations: Choosing the Right Tools.}
	\item Assignment 1:Basic Skills I: What you’ll have to know. 
	\vspace{1mm}
\end{itemize}
\end{minipage} \\
\hline

Week 2 & \begin{minipage}{.85\textwidth}
\begin{itemize} \itemsep-0.4em
	\vspace{1mm}
	\item 8/31 - Lecture; 9/2 and 9/4 - Hands On Learning
	\item Data Applications in  Decisionmaking: A History II  \item Reading: \textit{Can Selfishness Save the Environment?  The Atlantic Monthly, Available online}
	\item Assignment 1:Basic Skills II: What you’ll have to know. 
	\vspace{1mm}
\end{itemize}
\end{minipage} \\
\hline

Week 3 & \begin{minipage}{.85\textwidth}
\begin{itemize} \itemsep-0.4em
	\vspace{1mm}
	\item 9/7 - Lecture; 9/9 and 9/11 - Hands On Learning
	\item Challenge I: US Army Corps  \item Reading: \textit{Temporal and spatial changes in social vulnerability to natural hazards - Susan Cutter}
	\item Assignment 2:Working with Data: Figuring out what’s wrong.
	\item \textbf{"Basic Skills" Lab due by midnight on 9/11!}
	\vspace{1mm}
\end{itemize}
\end{minipage} \\
\hline

Week 4 & \begin{minipage}{.85\textwidth}
\begin{itemize} \itemsep-0.4em
	\vspace{1mm}
	\item 9/14 - Lecture; 9/16 and 9/18 - Hands On Learning
	\item Challenge I: Getting the Data  \item Reading: \textit{Vörösmarty, C., P. McIntyre, et al. (2010). "Global threats to human water security and river biodiversity." Nature 467(7315): 555-561.}
	\item Assignment 2: Working with Data: Doing it right.
	\vspace{1mm}
\end{itemize}
\end{minipage} \\
\hline

Week 5 & \begin{minipage}{.85\textwidth}
\begin{itemize} \itemsep-0.4em
	\vspace{1mm}
	\item 9/21 - Lecture; 9/23 and 9/25 - Hands On Learning
	\item Challenge I: Understanding the Data \item Reading: \textit{Chapter 15 in: Decision Analysis for Management Judgment: Second Edition, Goodwin and Wright, The Analytical Hierarchy Process.}
	\item Mid-Term Assignment [3]: Implementing the AHP –Working with Decisionmakers to make sense of data.
	\item \textbf{Assignment 2 (Working with Data) due by midnight on 9/25!}
	\vspace{1mm}
\end{itemize}
\end{minipage} \\
\hline


Week 6 & \begin{minipage}{.85\textwidth}
\begin{itemize} \itemsep-0.4em
	\vspace{1mm}
	\item 9/28 - Lecture; 9/30 and 10/2 - Hands On Learning
	\item Challenge I: Analyzing the Data 
	\item Reading: \textit{Clark, G., S. Moser, et al. (1998). "Assessing the vulnerability of coastal communities to extreme  storms: the case of Revere, MA., USA." Mitigation and Adaptation 
Strategies for Global Change 3(1): 59-82}
	\item Mid-Term Assignment [3]: Data Reduction: Drawing Inferences from “lots of data”
	\vspace{1mm}
\end{itemize}
\end{minipage} \\
\hline

Week 7 & \begin{minipage}{.85\textwidth}
\begin{itemize} \itemsep-0.4em
	\vspace{1mm}
	\item 10/5 - Lecture; 10/7 and 10/9 - Hands On Learning
	\item Challenge I: Making the Case
	\item Reading: \textit{Visualizing ourselves… with
crowd-sourced data – Aaron Koblin}
	\item No Assignment this week; preperation for mid-term (assignment 3) presentations.
	\vspace{1mm}
\end{itemize}
\end{minipage} \\
\hline

Week 8 & \begin{minipage}{.85\textwidth}
\begin{itemize} \itemsep-0.4em
	\vspace{1mm}
	\item 10/12 - No Class (Fall Break); 10/14 and 10/16 - Presentations.
	\item \textbf{Midterm Presentations and accompanying report (i.e. Assignment 3) due by midnight on Sunday, 10/11!}
	\item Reading: \textit{NA}
	\item NA
	\vspace{1mm}
\end{itemize}
\end{minipage} \\
\hline

\end{tabular} 
\end{table}

\begin{table}[h!]
\small % The size of the table text can be changed depending on content. Remove if desired.
\begin{tabular}{ | c | c | }
\hline
\textbf{Week} & \textbf{Content} \\
\hline

Week 9 & \begin{minipage}{.85\textwidth}
\begin{itemize} \itemsep-0.4em
	\vspace{1mm}
	\item 10/19 - Lecture; 10/21 and 10/23 - Hands On Learning.
	Selecting a Challenge and Forming Groups
	\item Reading: \textit{Student-Selected}
	\item Finding Data: Where do I start looking?
	\vspace{1mm}
\end{itemize}
\end{minipage} \\
\hline

Week 10 & \begin{minipage}{.85\textwidth}
\begin{itemize} \itemsep-0.4em
	\vspace{1mm}
	\item 10/26 - \textbf{Guest Lecture}; 10/28 and 10/30 - Hands On Learning.
	\item Challenge II: Getting the Data (Ready)
	\item Reading: \textit{Student-Selected}
	\item Working with Data: Cleaning up!
	\item \textbf{Reading summaries of at least one academic article on your topic of intereste due by midnight on Friday, 10/30 (reading assignment 1)!}
	\vspace{1mm}
\end{itemize}
\end{minipage} \\
\hline

Week 11 & \begin{minipage}{.85\textwidth}
\begin{itemize} \itemsep-0.4em
	\vspace{1mm}
	\item 11/2 - \textbf{Guest Lecture}; 11/4 and 11/6 - Hands On Learning.
	\item Challenge II: Getting the Data (Ready)
	\item Reading: \textit{Student-Selected}
	\item Working with Data: Cleaning up!
	\item \textbf{Dataset Due for your topic of interest by midnight on 11/6! (Counts as a lab assignment!)}
	\vspace{1mm}
\end{itemize}
\end{minipage} \\
\hline

Week 12 & \begin{minipage}{.85\textwidth}
\begin{itemize} \itemsep-0.4em
	\vspace{1mm}
	\item 11/9 -\textbf{ Guest Lecture}; 11/11 and 11/13 - Hands On Learning.
	\item Challenge II: Getting the Data Part 2
	\item Reading: \textit{Student-Selected}
	\item Working with Data: Getting ready for Analyses
	\item \textbf{Reading summaries of at least one academic article on your topic of interest due by midnight on Friday, 11/13 (reading assignment 2)!}
	\vspace{1mm}
\end{itemize}
\end{minipage} \\
\hline

Week 13 & \begin{minipage}{.85\textwidth}
\begin{itemize} \itemsep-0.4em
	\vspace{1mm}
	\item 11/16 - \textbf{Guest Lecture}; 11/18 and 11/20 - Hands On Learning.
	\item Challenge II: Analyzing the Data
	\item Reading: \textit{Student-Selected}
	\item Working with Data: Running the Analysis
	\item \textbf{First draft powerpoint presentation due (counts as a lab assignment!}
	\vspace{1mm}
\end{itemize}
\end{minipage} \\
\hline

Week 14 & \begin{minipage}{.85\textwidth}
\begin{itemize} \itemsep-0.4em
	\vspace{1mm}
	\item 11/23 - Hands On Learning; 11/25 and 11/27 - No Class (Thanksgiving)
	\item Challenge II: Analyzing the Data
	\item Working with Data: Building Visualizations
	\item \textbf{Reading summaries of at least one academic article on your topic of interest due by midnight on Monday, 11/23 (reading assignment 3)!}
	\vspace{1mm}
\end{itemize}
\end{minipage} \\
\hline

Week 15 & \begin{minipage}{.85\textwidth}
\begin{itemize} \itemsep-0.4em
	\vspace{1mm}
	\item 11/30, 12/2, 12/4 - Final Presentations
	\vspace{1mm}
\end{itemize}
\end{minipage} \\
\hline

\end{tabular} 
\end{table}

\end{document}



